\clearpage\thispagestyle{empty}
\pdfbookmark[0]{Sources for Calculation of Passover (1802)}{osterfest1802-sources}
\begin{center}\textbf{Sources for Calculation of Passover (1802)}\end{center}
\noindent This AI translation follows Gauss’s \textit{Berechnung des jüdischen Osterfestes} in \textit{Werke}, vol.~VI (Göttingen, 1874), pp.~80--81. The first printing appeared in Franz Xaver von Zach’s \textit{Monatliche Correspondenz zur Beförderung der Erd- und Himmels-Kunde}, vol.~5, May 1802, article~XLVI, pp.~435--437, under the name Dr.~Gauss in Braunschweig. \href{https://www.digitale-sammlungen.de/de/view/bsb10538597?page=447}{First printing}. \href{https://gdz.sub.uni-goettingen.de/id/PPN236007467}{Collected German text}.
\par\smallskip\noindent The title and text use \textit{Passover} for Gauss’s \textit{jüdisches Osterfest}. The English rendering retains the calendar terms and numerical forms in the collected edition, including the example’s extended March dates. For modern German usage see Barbara Staudinger, \href{https://www.jmw.at/news/die_pessach-formel}{\textit{Die Pessach-Formel}}, Jüdisches Museum Wien, 27 March 2025.
\par\smallskip\noindent Translation and typesetting produced by AI.
\clearpage
